\chapter{INTRODUCTION}

\section{Background}
In contemporary tertiary technical education, undergraduate engineering graduates face an increasingly demanding recruitment landscape. The software engineering industry has evolved from conventional algorithmic paper testing to multifaceted candidate screening pipelines. Today, premier technology companies, unicorns, and high-growth startups assess candidates based on a tripartite composite profile:
\begin{enumerate}
    \item \textbf{Practical Version-Controlled Development:} Observable evidence of real-world software creation through platforms such as GitHub, encompassing commit velocity, code modularity, repository licensing, and open-source contributions.
    \item \textbf{Problem-Solving and Algorithmic Rigor:} Competitive programming mastery demonstrated through platforms like LeetCode and HackerRank, emphasizing time-space algorithmic efficiency.
    \item \textbf{Industry-Aligned Curriculum Vitae (ATS Compatibility):} Professional resumes formatted to clear Automated Tracking Systems (ATS) that parse domain-specific keywords and impact metrics.
\end{enumerate}

Despite the proliferation of isolated tools for each of these dimensions, engineering students frequently struggle to gauge their holistic employability. This fragmentation creates a systemic diagnostic disconnect in undergraduate career readiness.

\section{Domain Overview}
The project operates at the confluence of \textbf{Automated Talent Assessment}, \textbf{Educational Data Mining (EDM)}, and \textbf{Applied Generative Artificial Intelligence (GenAI)}. Within university campus recruitment ecosystems, placement cells traditionally administer generic aptitude examinations and manual resume reviews. Such approaches suffer from distinct bottlenecks:
\begin{itemize}
    \item \textbf{Scalability Constraints:} University placement offices cater to thousands of graduating seniors with limited human evaluators, rendering individualized feedback impossible.
    \item \textbf{Static vs. Dynamic Competency:} Conventional assessment portals evaluate candidates through static multiple-choice questionnaires that fail to mirror realistic engineering environments.
    \item \textbf{Lack of Actionable Roadmaps:} When candidates fail screening rounds, existing systems provide binary rejection notifications without diagnosing conceptual deficiencies or prescribing adaptive remedial learning paths.
\end{itemize}

\section{Problem Statement}
Contemporary technical job aspirants lack a unified, intelligent platform capable of synthesizing developer telemetry, algorithmic proficiency, ATS resume parsing, and secure diagnostic assessment into a coherent employability index. Existing evaluation mechanisms exhibit three critical problems:
\begin{enumerate}
    \item \textbf{Siloed Portfolios:} No existing system automatically ingests and cross-analyzes a student's GitHub activity and LeetCode performance within a single diagnostic dashboard.
    \item \textbf{Superficial Resume Feedback:} Job applicants receive generic ATS percentage scores without contextual phrase-level semantic recommendations or role-specific gap analysis.
    \item \textbf{Vulnerable Mock Assessments:} Standard web-based mock examinations are plagued by unmonitored tab-switching, unauthorized clipboard copy-pasting, and AI-assisted screen capture cheating.
\end{enumerate}

\section{Motivation}
The inspiration behind developing \textbf{PlacementPulse} stems from direct observations of the hurdles experienced by computer science students during campus placement seasons. Capable students are frequently disqualified at initial screening stages due to low ATS match scores, unverified GitHub hygiene, or inability to handle timed proctored assessments. A technological solution combining automated API telemetry ingestion, browser-level proctoring heuristics, and multi-model LLM orchestration democratizes high-tier placement preparation, ensuring that every student receives equitable, objective, and personalized career mentoring.

\section{Objectives}
The primary objective of the PlacementPulse project is to design, develop, and deploy an end-to-end, full-stack placement readiness ecosystem. The specific sub-objectives are:
\begin{enumerate}
    \item \textbf{Unified Profile Integration:} Build automated ingestion connectors for GitHub REST/GraphQL APIs and LeetCode GraphQL endpoints to evaluate repository quality, language diversity, and algorithmic difficulty breakdowns.
    \item \textbf{Intelligent ATS Resume Parser:} Develop a server-side PDF extraction engine utilizing keyword clustering to compute role-tailored ATS match scores and missing skills.
    \item \textbf{Secure Proctored Assessment Engine:} Engineer a 52-question randomized diagnostic exam portal featuring DOM canvas text-rendering, clipboard disabling, and window blur tracking.
    \item \textbf{Autonomous Multi-Model AI Mentor:} Integrate an LLM-powered pedagogical agent utilizing an automated multi-model fallback chain (Llama-3.3-70B $\rightarrow$ DeepSeek-Chat $\rightarrow$ Qwen-2.5-72B) to guarantee zero downtime during conversational tutoring.
    \item \textbf{Adaptive Learning Roadmap Generation:} Formulate an algorithm that synthesizes diagnostic gaps to generate structured, milestone-driven technical roadmaps.
\end{enumerate}

\section{Scope of the Project}
The functional scope encompasses user authentication with email OTP verification, public developer portfolio ingestion, in-memory PDF ATS auditing, a timed 52-question canvas-secured examination across core Computer Science domains, and multi-model AI mentoring. External biometric hardware surveillance and continuous video/audio streaming are excluded to preserve low bandwidth overhead and student privacy.

\section{Target Users}
PlacementPulse is tailored for pre-final and final year engineering undergraduates preparing for campus drives, early-career software engineering aspirants seeking portfolio audits, and university placement coordinators requiring objective cohort diagnostic analytics.
